\documentclass{beamer}
\usetheme{Madrid}
\usecolortheme{default}

% Define OSU orange color
\definecolor{osaorange}{RGB}{220,115,38}

% Set colors
\setbeamercolor{structure}{fg=osaorange}
\setbeamercolor{title}{fg=white,bg=osaorange}
\setbeamercolor{frametitle}{fg=white,bg=osaorange}

\usepackage{graphicx}
\usepackage{amsmath}
\usepackage{booktabs}
\usepackage{hyperref}
\usepackage{tikz}
\usetikzlibrary{shapes,arrows}

\title{H524 Introduction to Biostatistics}
\subtitle{Week 9: Correlation and Simple Linear Regression}
\author{Dr. John Molitor}
\institute{Oregon State University\\College of Health}
\date{Fall 2025}

\begin{document}

\frame{\titlepage}

\begin{frame}
\frametitle{Outline}
\tableofcontents
\end{frame}

\section{Review and Introduction}

\begin{frame}
\frametitle{What We've Learned So Far}
\textbf{Weeks 1-8 covered:}
\begin{itemize}
\item Descriptive statistics and data visualization
\item Probability and diagnostic testing
\item Sampling distributions and Central Limit Theorem
\item Confidence intervals
\item Hypothesis testing (one-sample, two-sample, paired)
\item ANOVA (comparing multiple groups)
\item Categorical data analysis (chi-square tests)
\end{itemize}

\vspace{0.5cm}
\textbf{This Week:} Relationships between continuous variables
\begin{itemize}
\item How do two variables relate to each other?
\item Can we predict one variable from another?
\end{itemize}

\vspace{0.3cm}
\textbf{Reading:} Pagano \& Gauvreau Chapters 17 and 18
\end{frame}

\begin{frame}
\frametitle{New Questions We Can Answer}
\textbf{Previous weeks:} Comparisons between groups
\begin{itemize}
\item Is mean cholesterol different between treatment groups?
\item Is smoking status associated with lung cancer?
\end{itemize}

\vspace{0.5cm}
\textbf{This week:} Relationships between continuous variables
\begin{itemize}
\item How is systolic blood pressure related to age?
\item Can we predict birth weight from gestational age?
\item Is there a relationship between BMI and cholesterol?
\item How strong is the association between height and weight?
\end{itemize}

\vspace{0.5cm}
\textbf{Key concepts:}
\begin{itemize}
\item \textbf{Correlation:} Measure strength and direction of linear relationship
\item \textbf{Regression:} Predict one variable from another
\end{itemize}
\end{frame}

\section{Correlation (Chapter 17)}

\begin{frame}
\frametitle{What is Correlation?}
\textbf{Correlation:} A measure of the linear association between two continuous variables

\vspace{0.5cm}
\textbf{Key characteristics:}
\begin{itemize}
\item Measures strength and direction of relationship
\item Ranges from $-1$ to $+1$
\item $r = 0$: No linear relationship
\item $r = +1$: Perfect positive linear relationship
\item $r = -1$: Perfect negative linear relationship
\item $|r| > 0.7$: Strong correlation
\item $|r| = 0.3$ to $0.7$: Moderate correlation
\item $|r| < 0.3$: Weak correlation
\end{itemize}

\vspace{0.3cm}
\textbf{Important:} Correlation measures LINEAR relationships only
\begin{itemize}
\item Two variables can be strongly related but have $r \approx 0$ if relationship is nonlinear
\end{itemize}
\end{frame}

\begin{frame}
\frametitle{Visualizing Correlation}
\textbf{Scatter plots reveal the relationship:}

\vspace{0.3cm}
\begin{itemize}
\item \textbf{Positive correlation:} As X increases, Y increases
  \begin{itemize}
  \item Example: Height and weight
  \item Points trend upward from left to right
  \end{itemize}

\vspace{0.3cm}
\item \textbf{Negative correlation:} As X increases, Y decreases
  \begin{itemize}
  \item Example: Exercise frequency and resting heart rate
  \item Points trend downward from left to right
  \end{itemize}

\vspace{0.3cm}
\item \textbf{No correlation:} No clear pattern
  \begin{itemize}
  \item Example: Shoe size and IQ
  \item Points scattered randomly
  \end{itemize}
\end{itemize}

\vspace{0.3cm}
\textbf{Always plot your data first!} Summary statistics can mislead without visualization.
\end{frame}

\begin{frame}
\frametitle{Pearson Correlation Coefficient}
\textbf{Most common correlation measure: Pearson's $r$}

\vspace{0.3cm}
\textbf{Formula:}
$$r = \frac{\sum_{i=1}^{n}(x_i - \bar{x})(y_i - \bar{y})}{\sqrt{\sum_{i=1}^{n}(x_i - \bar{x})^2} \sqrt{\sum_{i=1}^{n}(y_i - \bar{y})^2}}$$

Alternative formula:
$$r = \frac{1}{n-1} \sum_{i=1}^{n} \left(\frac{x_i - \bar{x}}{s_x}\right) \left(\frac{y_i - \bar{y}}{s_y}\right)$$

\vspace{0.3cm}
\textbf{Interpretation:}
\begin{itemize}
\item $r$ measures how consistently X and Y vary together
\item Standardized measure (unitless)
\item Symmetric: $r_{xy} = r_{yx}$
\end{itemize}

\vspace{0.3cm}
\textbf{In R:} \texttt{cor(x, y)} or \texttt{cor.test(x, y)}
\end{frame}

\begin{frame}
\frametitle{Example: Age and Systolic Blood Pressure}
\textbf{Research question:} Is blood pressure related to age?

\vspace{0.3cm}
\textbf{Data:} 15 adults aged 30-70 years

\begin{center}
\begin{tabular}{ccc|ccc}
\toprule
Subject & Age & SBP & Subject & Age & SBP \\
\midrule
1 & 35 & 118 & 9 & 52 & 138 \\
2 & 38 & 122 & 10 & 55 & 140 \\
3 & 40 & 125 & 11 & 58 & 145 \\
4 & 42 & 128 & 12 & 60 & 148 \\
5 & 45 & 130 & 13 & 62 & 150 \\
6 & 47 & 132 & 14 & 65 & 155 \\
7 & 48 & 135 & 15 & 68 & 158 \\
8 & 50 & 136 & & & \\
\bottomrule
\end{tabular}
\end{center}

\vspace{0.3cm}
\textbf{Pearson correlation:} $r = 0.98$ (very strong positive)

\textbf{Interpretation:} Age and systolic BP are strongly positively correlated
\end{frame}

\begin{frame}
\frametitle{Testing Correlation Significance}
\textbf{Question:} Is the observed correlation statistically significant, or could it be due to chance?

\vspace{0.3cm}
\textbf{Hypotheses:}
\begin{itemize}
\item $H_0: \rho = 0$ (population correlation is zero)
\item $H_A: \rho \neq 0$ (population correlation is not zero)
\end{itemize}

\vspace{0.3cm}
\textbf{Test statistic:}
$$t = \frac{r\sqrt{n-2}}{\sqrt{1-r^2}} \sim t_{n-2}$$

\vspace{0.3cm}
\textbf{Example: Age and BP}
\begin{itemize}
\item $r = 0.98$, $n = 15$
\item $t = \frac{0.98\sqrt{13}}{\sqrt{1-0.98^2}} = \frac{3.535}{0.199} = 17.76$
\item $df = 13$, $p < 0.001$
\item \textbf{Conclusion:} Very strong evidence of positive correlation
\end{itemize}

\vspace{0.3cm}
\textbf{In R:} \texttt{cor.test(age, sbp)} provides $r$, $t$, $p$-value, and CI
\end{frame}

\begin{frame}
\frametitle{Assumptions for Pearson Correlation}
\textbf{For valid inference, we assume:}

\begin{enumerate}
\item \textbf{Bivariate normality:} (X, Y) follow a bivariate normal distribution
  \begin{itemize}
  \item Each variable approximately normal
  \item Joint distribution is elliptical
  \end{itemize}

\item \textbf{Linear relationship:} Association is linear
  \begin{itemize}
  \item Check with scatter plot
  \item Pearson $r$ may miss nonlinear patterns
  \end{itemize}

\item \textbf{No extreme outliers:} Outliers strongly influence $r$
  \begin{itemize}
  \item Single outlier can create or destroy correlation
  \item Always examine scatter plot
  \end{itemize}

\item \textbf{Random sample:} Observations independently sampled
\end{enumerate}

\vspace{0.3cm}
\textbf{If assumptions violated:} Consider Spearman's rank correlation (more robust)
\end{frame}

\begin{frame}
\frametitle{Spearman's Rank Correlation}
\textbf{Alternative to Pearson:} Based on ranks, not raw values

\vspace{0.3cm}
\textbf{When to use Spearman's $r_s$:}
\begin{itemize}
\item Relationship is monotonic but not linear
\item Data contain outliers
\item Data are ordinal (ranks)
\item Distribution is highly skewed
\end{itemize}

\vspace{0.3cm}
\textbf{Method:}
\begin{enumerate}
\item Rank X values: $1, 2, 3, \ldots, n$
\item Rank Y values: $1, 2, 3, \ldots, n$
\item Calculate Pearson correlation on ranks
\end{enumerate}

\vspace{0.3cm}
\textbf{Properties:}
\begin{itemize}
\item Ranges from $-1$ to $+1$
\item Less sensitive to outliers than Pearson $r$
\item Detects monotonic (not just linear) relationships
\end{itemize}

\vspace{0.3cm}
\textbf{In R:} \texttt{cor(x, y, method="spearman")}
\end{frame}

\begin{frame}
\frametitle{Correlation vs. Causation}
\textbf{CRITICAL WARNING:}

\vspace{0.3cm}
\begin{center}
\Large{\textbf{Correlation does NOT imply causation!}}
\end{center}

\vspace{0.5cm}
\textbf{Examples of spurious correlations:}
\begin{itemize}
\item Ice cream sales and drowning deaths (both increase in summer)
\item Number of fire trucks and damage from fires (both higher for big fires)
\item Shoe size and reading ability in children (both increase with age)
\end{itemize}

\vspace{0.5cm}
\textbf{Possible explanations for correlation:}
\begin{enumerate}
\item X causes Y
\item Y causes X
\item Both X and Y caused by third variable (confounding)
\item Coincidence
\end{enumerate}

\vspace{0.3cm}
\textbf{To establish causation:} Need randomized controlled experiments, not just observational correlations
\end{frame}

\begin{frame}
\frametitle{Important Properties of Correlation}
\textbf{What correlation tells us:}
\begin{itemize}
\item Strength of linear relationship
\item Direction (positive or negative)
\item Statistical significance
\end{itemize}

\vspace{0.3cm}
\textbf{What correlation does NOT tell us:}
\begin{itemize}
\item Causation (as discussed)
\item The form of the relationship (linear? exponential?)
\item How to predict Y from X
\end{itemize}

\vspace{0.5cm}
\textbf{Key limitations:}
\begin{itemize}
\item Sensitive to outliers
\item Only measures linear relationships
\item Same $r$ can arise from very different scatter plots
\item Restricted range reduces correlation
\end{itemize}

\vspace{0.3cm}
\textbf{Coefficient of determination:} $r^2$
\begin{itemize}
\item Proportion of variability in Y explained by X
\item Example: $r = 0.80$ means $r^2 = 0.64$ (64\% of variance explained)
\end{itemize}
\end{frame}

\section{Simple Linear Regression (Chapter 18)}

\begin{frame}
\frametitle{From Correlation to Regression}
\textbf{Correlation:} Measures strength of relationship (symmetric)

\vspace{0.3cm}
\textbf{Regression:} Models relationship to make predictions (asymmetric)

\vspace{0.5cm}
\textbf{Simple linear regression:}
\begin{itemize}
\item One response variable (Y): outcome we want to predict
\item One predictor variable (X): variable we use for prediction
\item Assumes linear relationship: $Y = \beta_0 + \beta_1 X + \epsilon$
\end{itemize}

\vspace{0.5cm}
\textbf{Goals of regression:}
\begin{enumerate}
\item Describe relationship between X and Y
\item Estimate the effect of X on Y
\item Predict Y for given values of X
\item Test hypotheses about the relationship
\end{enumerate}

\vspace{0.3cm}
\textbf{Example:} Predict birth weight (Y) from gestational age (X)
\end{frame}

\begin{frame}
\frametitle{The Simple Linear Regression Model}
\textbf{Population model:}
$$Y_i = \beta_0 + \beta_1 X_i + \epsilon_i$$

where:
\begin{itemize}
\item $Y_i$: Response variable for subject $i$
\item $X_i$: Predictor variable for subject $i$
\item $\beta_0$: Population intercept (value of Y when X = 0)
\item $\beta_1$: Population slope (change in Y per unit change in X)
\item $\epsilon_i$: Random error term, $\epsilon_i \sim N(0, \sigma^2)$
\end{itemize}

\vspace{0.5cm}
\textbf{Sample regression line (fitted model):}
$$\hat{Y}_i = b_0 + b_1 X_i$$

where:
\begin{itemize}
\item $\hat{Y}_i$: Predicted value of Y for subject $i$
\item $b_0$: Sample estimate of $\beta_0$
\item $b_1$: Sample estimate of $\beta_1$
\end{itemize}
\end{frame}

\begin{frame}
\frametitle{Interpreting Regression Coefficients}
\textbf{Intercept ($b_0$):}
\begin{itemize}
\item Predicted value of Y when X = 0
\item May not be meaningful if X = 0 is outside observed range
\item Example: Birth weight when gestational age = 0 weeks (nonsensical)
\end{itemize}

\vspace{0.5cm}
\textbf{Slope ($b_1$):}
\begin{itemize}
\item Change in Y for each 1-unit increase in X
\item Most important parameter
\item \textbf{Example:} $b_1 = 100$ g/week
  \begin{itemize}
  \item Each additional week of gestation increases birth weight by 100 g
  \item 3 weeks longer gestation $\Rightarrow$ 300 g heavier baby (on average)
  \end{itemize}
\end{itemize}

\vspace{0.5cm}
\textbf{Key point:} Slope quantifies the relationship
\begin{itemize}
\item $b_1 > 0$: Positive relationship
\item $b_1 < 0$: Negative relationship
\item $b_1 = 0$: No linear relationship
\end{itemize}
\end{frame}

\begin{frame}
\frametitle{Least Squares Estimation}
\textbf{How do we find the best-fitting line?}

\vspace{0.3cm}
\textbf{Method of least squares:} Minimize sum of squared residuals

\vspace{0.3cm}
\textbf{Residual:} $e_i = Y_i - \hat{Y}_i$ (observed $-$ predicted)

\vspace{0.3cm}
\textbf{Objective:} Minimize $\sum_{i=1}^{n} e_i^2 = \sum_{i=1}^{n} (Y_i - \hat{Y}_i)^2$

\vspace{0.5cm}
\textbf{Formulas for estimates:}
$$b_1 = \frac{\sum_{i=1}^{n}(X_i - \bar{X})(Y_i - \bar{Y})}{\sum_{i=1}^{n}(X_i - \bar{X})^2} = r \times \frac{s_Y}{s_X}$$

$$b_0 = \bar{Y} - b_1 \bar{X}$$

\vspace{0.3cm}
\textbf{Properties:}
\begin{itemize}
\item Regression line always passes through $(\bar{X}, \bar{Y})$
\item Unique solution (one best-fitting line)
\item Sum of residuals equals zero: $\sum e_i = 0$
\end{itemize}
\end{frame}

\begin{frame}
\frametitle{Example: Age and Blood Pressure Regression}
\textbf{Data:} Age (X) and systolic BP (Y) from 15 adults

\vspace{0.3cm}
\textbf{Summary statistics:}
\begin{itemize}
\item $\bar{X} = 50$ years, $s_X = 10.49$ years
\item $\bar{Y} = 138$ mmHg, $s_Y = 13.42$ mmHg
\item $r = 0.98$
\end{itemize}

\vspace{0.3cm}
\textbf{Calculate slope:}
$$b_1 = r \times \frac{s_Y}{s_X} = 0.98 \times \frac{13.42}{10.49} = 1.25 \text{ mmHg/year}$$

\textbf{Calculate intercept:}
$$b_0 = \bar{Y} - b_1 \bar{X} = 138 - 1.25(50) = 75.5 \text{ mmHg}$$

\vspace{0.3cm}
\textbf{Regression equation:}
$$\widehat{\text{SBP}} = 75.5 + 1.25 \times \text{Age}$$

\textbf{Interpretation:} For each additional year of age, systolic BP increases by 1.25 mmHg on average
\end{frame}

\begin{frame}
\frametitle{Making Predictions}
\textbf{Using the regression equation for prediction:}
$$\widehat{\text{SBP}} = 75.5 + 1.25 \times \text{Age}$$

\vspace{0.3cm}
\textbf{Example predictions:}
\begin{itemize}
\item Age 40: $\widehat{\text{SBP}} = 75.5 + 1.25(40) = 125.5$ mmHg
\item Age 55: $\widehat{\text{SBP}} = 75.5 + 1.25(55) = 144.3$ mmHg
\item Age 70: $\widehat{\text{SBP}} = 75.5 + 1.25(70) = 163.0$ mmHg
\end{itemize}

\vspace{0.5cm}
\textbf{Cautions about prediction:}
\begin{itemize}
\item \textbf{Interpolation:} Predictions within range of observed X (safe)
\item \textbf{Extrapolation:} Predictions outside range of observed X (risky!)
  \begin{itemize}
  \item Age 20: $\widehat{\text{SBP}} = 75.5 + 1.25(20) = 100.5$ mmHg (OK if reasonable)
  \item Age 100: $\widehat{\text{SBP}} = 75.5 + 1.25(100) = 200.5$ mmHg (unreliable!)
  \end{itemize}
\item Relationship may not hold outside observed range
\item Always check if prediction is in plausible range
\end{itemize}
\end{frame}

\begin{frame}
\frametitle{Residuals and Residual Plots}
\textbf{Residual:} Difference between observed and predicted
$$e_i = Y_i - \hat{Y}_i$$

\vspace{0.3cm}
\textbf{Why residuals matter:}
\begin{itemize}
\item Measure quality of fit
\item Check regression assumptions
\item Identify outliers and influential points
\end{itemize}

\vspace{0.5cm}
\textbf{Residual plot:} Plot $e_i$ vs. $\hat{Y}_i$ (or vs. $X_i$)

\vspace{0.3cm}
\textbf{What to look for:}
\begin{itemize}
\item \textbf{Good:} Random scatter around zero, constant spread
\item \textbf{Problem:} Curved pattern (nonlinearity)
\item \textbf{Problem:} Funnel shape (non-constant variance)
\item \textbf{Problem:} Extreme outliers
\end{itemize}

\vspace{0.3cm}
\textbf{Always examine residual plots!} They reveal model inadequacies
\end{frame}

\begin{frame}
\frametitle{Assessing Model Fit: $R^2$}
\textbf{Coefficient of determination:} $R^2$

\vspace{0.3cm}
$$R^2 = \frac{\text{SS}_{\text{regression}}}{\text{SS}_{\text{total}}} = 1 - \frac{\text{SS}_{\text{residual}}}{\text{SS}_{\text{total}}}$$

For simple linear regression: $R^2 = r^2$

\vspace{0.5cm}
\textbf{Interpretation:}
\begin{itemize}
\item Proportion of variability in Y explained by X
\item Ranges from 0 to 1 (or 0\% to 100\%)
\item $R^2 = 0$: X explains none of the variability in Y
\item $R^2 = 1$: X explains all of the variability in Y (perfect fit)
\end{itemize}

\vspace{0.3cm}
\textbf{Example: Age and BP}
\begin{itemize}
\item $r = 0.98$, so $R^2 = 0.98^2 = 0.96$ or 96\%
\item Age explains 96\% of the variability in systolic BP
\item Remaining 4\% due to other factors (genetics, diet, exercise, etc.)
\end{itemize}
\end{frame}

\begin{frame}
\frametitle{Regression Assumptions}
\textbf{For valid inference, we assume:}

\begin{enumerate}
\item \textbf{Linearity:} True relationship between X and Y is linear
  \begin{itemize}
  \item Check with scatter plot and residual plot
  \end{itemize}

\item \textbf{Independence:} Observations are independent
  \begin{itemize}
  \item Violated by clustered or repeated measures data
  \end{itemize}

\item \textbf{Normality:} Residuals are normally distributed
  \begin{itemize}
  \item Check with Q-Q plot or histogram of residuals
  \item Less critical for large samples (CLT)
  \end{itemize}

\item \textbf{Constant variance (homoscedasticity):} Spread of residuals is constant
  \begin{itemize}
  \item Check with residual plot
  \item No funnel shape
  \end{itemize}

\item \textbf{No influential outliers:} Extreme points can distort fit
  \begin{itemize}
  \item Check with residual plot and leverage diagnostics
  \end{itemize}
\end{enumerate}
\end{frame}

\begin{frame}
\frametitle{Inference for Regression Slope}
\textbf{Key question:} Is there a significant linear relationship between X and Y?

\vspace{0.3cm}
\textbf{Hypotheses:}
\begin{itemize}
\item $H_0: \beta_1 = 0$ (no linear relationship)
\item $H_A: \beta_1 \neq 0$ (linear relationship exists)
\end{itemize}

\vspace{0.3cm}
\textbf{Test statistic:}
$$t = \frac{b_1 - 0}{\text{SE}(b_1)} \sim t_{n-2}$$

where $\text{SE}(b_1) = \frac{s}{\sqrt{\sum(X_i - \bar{X})^2}}$ and $s = \sqrt{\frac{\sum e_i^2}{n-2}}$ is the residual standard error

\vspace{0.3cm}
\textbf{Confidence interval for $\beta_1$:}
$$b_1 \pm t_{n-2, \alpha/2} \times \text{SE}(b_1)$$

\vspace{0.3cm}
\textbf{Note:} In simple linear regression, testing $\beta_1 = 0$ is equivalent to testing $\rho = 0$
\end{frame}

\begin{frame}
\frametitle{Example: Testing Age-BP Relationship}
\textbf{Regression output from R:}

\begin{center}
\small
\begin{tabular}{lrrrr}
\toprule
& Estimate & Std. Error & t value & $p$-value \\
\midrule
Intercept & 75.5 & 4.23 & 17.85 & $<0.001$ \\
Age & 1.25 & 0.083 & 15.06 & $<0.001$ \\
\bottomrule
\end{tabular}
\end{center}

Residual standard error: 2.84 mmHg\\
$R^2 = 0.946$, Adjusted $R^2 = 0.942$

\vspace{0.3cm}
\textbf{Interpretation:}
\begin{itemize}
\item $b_1 = 1.25$ mmHg/year (95\% CI: 1.07 to 1.43)
\item $t = 15.06$ with $df = 13$, $p < 0.001$
\item Very strong evidence that age is associated with systolic BP
\item For each year increase in age, SBP increases by 1.25 mmHg (95\% CI: 1.07 to 1.43)
\end{itemize}

\vspace{0.3cm}
\textbf{In R:} \texttt{summary(lm(sbp \textasciitilde\ age))}
\end{frame}

\begin{frame}
\frametitle{ANOVA for Regression}
\textbf{Analysis of Variance decomposes total variability:}

$$\text{SS}_{\text{total}} = \text{SS}_{\text{regression}} + \text{SS}_{\text{residual}}$$

\vspace{0.3cm}
\begin{center}
\begin{tabular}{lrrrr}
\toprule
Source & df & SS & MS & F \\
\midrule
Regression & 1 & 1823.7 & 1823.7 & 226.8 \\
Residual & 13 & 104.5 & 8.04 & \\
\midrule
Total & 14 & 1928.2 & & \\
\bottomrule
\end{tabular}
\end{center}

\vspace{0.3cm}
\textbf{F-test:}
$$F = \frac{\text{MS}_{\text{regression}}}{\text{MS}_{\text{residual}}} = \frac{1823.7}{8.04} = 226.8$$

$p < 0.001$: Strong evidence that the regression model is significant

\vspace{0.3cm}
\textbf{Note:} For simple linear regression, $F = t^2$ and both tests are equivalent
\end{frame}

\begin{frame}
\frametitle{Prediction Intervals}
\textbf{Two types of intervals for predictions:}

\vspace{0.3cm}
\textbf{1. Confidence interval for mean response:}
\begin{itemize}
\item Estimates the mean Y for a given X
\item Narrower interval
\item ``What is the average BP for all 55-year-olds?''
\end{itemize}

\vspace{0.3cm}
\textbf{2. Prediction interval for individual response:}
\begin{itemize}
\item Predicts Y for a specific individual at given X
\item Wider interval (accounts for individual variation)
\item ``What BP will this specific 55-year-old have?''
\end{itemize}

\vspace{0.5cm}
\textbf{Both intervals:}
\begin{itemize}
\item Are narrowest at $\bar{X}$
\item Widen as we move away from $\bar{X}$
\item Should not be used for extrapolation
\end{itemize}

\vspace{0.3cm}
\textbf{In R:} \texttt{predict(model, newdata, interval="confidence")} or \texttt{interval="prediction"}
\end{frame}

\begin{frame}
\frametitle{Example: Prediction Intervals}
\textbf{For a 55-year-old:}

\vspace{0.3cm}
Point prediction: $\widehat{\text{SBP}} = 75.5 + 1.25(55) = 144.3$ mmHg

\vspace{0.5cm}
\textbf{95\% Confidence interval for mean SBP:}
\begin{itemize}
\item Estimates average SBP for all 55-year-olds
\item CI: (142.8, 145.8) mmHg
\item Relatively narrow (2-3 mmHg)
\end{itemize}

\vspace{0.5cm}
\textbf{95\% Prediction interval for individual SBP:}
\begin{itemize}
\item Predicts SBP for one specific 55-year-old
\item PI: (138.0, 150.6) mmHg
\item Much wider (12-13 mmHg)
\item Accounts for individual variation around the mean
\end{itemize}

\vspace{0.3cm}
\textbf{Interpretation:} We're 95\% confident this individual's SBP will be between 138 and 151 mmHg
\end{frame}

\begin{frame}
\frametitle{Transformations in Regression}
\textbf{When assumptions are violated:} Consider transforming variables

\vspace{0.3cm}
\textbf{Common transformations:}
\begin{itemize}
\item \textbf{Log transformation:} For right-skewed data or multiplicative relationships
  \begin{itemize}
  \item $\log(Y)$ vs. $X$: Exponential growth/decay
  \item $Y$ vs. $\log(X)$: Logarithmic relationship
  \item $\log(Y)$ vs. $\log(X)$: Power relationship
  \end{itemize}

\item \textbf{Square root:} For count data or moderate skewness
  \begin{itemize}
  \item $\sqrt{Y}$ vs. $X$
  \end{itemize}

\item \textbf{Reciprocal:} For hyperbolic relationships
  \begin{itemize}
  \item $1/Y$ vs. $X$
  \end{itemize}
\end{itemize}

\vspace{0.3cm}
\textbf{Benefits:}
\begin{itemize}
\item Linearize relationship
\item Stabilize variance
\item Improve normality of residuals
\end{itemize}

\vspace{0.3cm}
\textbf{Caution:} Interpretation changes with transformed variables
\end{frame}

\section{Summary and Examples}

\begin{frame}
\frametitle{Correlation vs. Regression Summary}
\begin{center}
\begin{tabular}{p{4cm}p{4cm}}
\toprule
\textbf{Correlation} & \textbf{Regression} \\
\midrule
Symmetric & Asymmetric \\
Measures strength & Models relationship \\
No prediction & Makes predictions \\
$r$ ranges $-1$ to $+1$ & Slope can be any value \\
Unitless & Has units \\
$r_{xy} = r_{yx}$ & Different X vs. Y roles \\
Single number & Equation \\
\bottomrule
\end{tabular}
\end{center}

\vspace{0.3cm}
\textbf{When to use each:}
\begin{itemize}
\item \textbf{Correlation:} To measure and test strength of association
\item \textbf{Regression:} To predict, model relationships, estimate effects
\end{itemize}

\vspace{0.3cm}
\textbf{Relationship:} In simple linear regression, $t$-test for $\beta_1 = 0$ is equivalent to $t$-test for $\rho = 0$
\end{frame}

\begin{frame}[fragile]
\frametitle{R Code: Complete Analysis}
\footnotesize
\begin{verbatim}
# Data
age <- c(35, 38, 40, 42, 45, 47, 48, 50, 52, 55,
         58, 60, 62, 65, 68)
sbp <- c(118, 122, 125, 128, 130, 132, 135, 136, 138,
         140, 145, 148, 150, 155, 158)

# 1. Correlation
cor.test(age, sbp)  # r = 0.98, p < 0.001

# 2. Scatter plot
plot(age, sbp, main = "Age vs. SBP",
     xlab = "Age (years)", ylab = "SBP (mmHg)",
     pch = 19, col = "blue")

# 3. Regression
model <- lm(sbp ~ age)
summary(model)      # Detailed output
abline(model, col = "red", lwd = 2)  # Add line to plot

# 4. Diagnostics
par(mfrow = c(2, 2))
plot(model)         # All diagnostic plots

# 5. Predictions
predict(model, newdata = data.frame(age = 55),
        interval = "prediction")
\end{verbatim}
\end{frame}

\begin{frame}
\frametitle{Practical Example: Birth Weight Study}
\textbf{Research question:} How does gestational age affect birth weight?

\vspace{0.3cm}
\textbf{Data:} 50 newborns, gestational age 32-42 weeks

\vspace{0.3cm}
\textbf{Analysis:}
\begin{itemize}
\item Scatter plot shows strong positive linear relationship
\item Pearson correlation: $r = 0.87$ ($p < 0.001$)
\item Regression: $\widehat{\text{Weight}} = -2800 + 150 \times \text{Weeks}$
  \begin{itemize}
  \item Intercept: $-2800$ g (not meaningful; 0 weeks is impossible)
  \item Slope: 150 g/week (95\% CI: 130 to 170)
  \item $R^2 = 0.76$ (76\% of variability explained)
  \end{itemize}
\end{itemize}

\vspace{0.3cm}
\textbf{Interpretation:}
\begin{itemize}
\item Each additional week of gestation increases birth weight by 150 g on average
\item Strong evidence of positive linear relationship ($p < 0.001$)
\item Can predict birth weight from gestational age with reasonable accuracy
\end{itemize}
\end{frame}

\begin{frame}
\frametitle{Practical Example: Exercise and Heart Rate}
\textbf{Research question:} Is exercise frequency associated with resting heart rate?

\vspace{0.3cm}
\textbf{Data:} 30 adults, exercise sessions per week (0-7), resting HR (bpm)

\vspace{0.3cm}
\textbf{Analysis:}
\begin{itemize}
\item Scatter plot shows negative linear trend
\item Pearson correlation: $r = -0.68$ ($p < 0.001$)
\item Regression: $\widehat{\text{HR}} = 78 - 2.5 \times \text{Sessions}$
  \begin{itemize}
  \item Intercept: 78 bpm (resting HR with no exercise)
  \item Slope: $-2.5$ bpm/session (95\% CI: $-3.8$ to $-1.2$)
  \item $R^2 = 0.46$ (46\% explained)
  \end{itemize}
\end{itemize}

\vspace{0.3cm}
\textbf{Interpretation:}
\begin{itemize}
\item Each additional exercise session per week reduces resting HR by 2.5 bpm
\item Moderate negative association (still significant)
\item Other factors also influence resting HR (54\% of variance unexplained)
\end{itemize}
\end{frame}

\begin{frame}
\frametitle{Common Pitfalls in Regression}
\textbf{Mistakes to avoid:}

\begin{enumerate}
\item \textbf{Ignoring assumptions}
  \begin{itemize}
  \item Always check residual plots
  \item Don't force linear model on nonlinear data
  \end{itemize}

\item \textbf{Extrapolation}
  \begin{itemize}
  \item Don't predict outside range of observed X
  \item Relationship may not hold in unobserved region
  \end{itemize}

\item \textbf{Causation claims}
  \begin{itemize}
  \item Regression shows association, not causation
  \item Need experiments for causal inference
  \end{itemize}

\item \textbf{Ignoring outliers}
  \begin{itemize}
  \item Single outlier can drastically change regression line
  \item Investigate influential points
  \end{itemize}

\item \textbf{Over-interpreting $R^2$}
  \begin{itemize}
  \item Low $R^2$ doesn't mean relationship is unimportant
  \item Clinical significance $\neq$ statistical significance
  \end{itemize}
\end{enumerate}
\end{frame}

\begin{frame}
\frametitle{Summary}
\textbf{Today we covered:}
\begin{itemize}
\item \textbf{Chapter 17: Correlation}
  \begin{itemize}
  \item Pearson and Spearman correlation
  \item Testing significance
  \item Correlation vs. causation
  \end{itemize}

\item \textbf{Chapter 18: Simple Linear Regression}
  \begin{itemize}
  \item Least squares estimation
  \item Interpreting slope and intercept
  \item Inference for regression coefficients
  \item Prediction and prediction intervals
  \item Checking assumptions with residual plots
  \end{itemize}
\end{itemize}

\vspace{0.5cm}
\textbf{Key formulas:}
\begin{itemize}
\item Regression line: $\hat{Y} = b_0 + b_1 X$
\item Slope: $b_1 = r \times s_Y / s_X$
\item $R^2 = r^2$ (proportion of variance explained)
\end{itemize}

\vspace{0.3cm}
\textbf{Next Week:} Multiple regression (more than one predictor)
\end{frame}

\begin{frame}
\frametitle{Key Takeaways}
\begin{enumerate}
\item \textbf{Always visualize first:} Scatter plots reveal patterns

\item \textbf{Correlation measures association:} But not causation

\item \textbf{Regression quantifies relationships:} Allows prediction

\item \textbf{Check assumptions:} Use residual plots

\item \textbf{Interpret slope in context:} What does 1-unit change in X mean?

\item \textbf{$R^2$ shows explained variance:} But low $R^2$ can still be important

\item \textbf{Don't extrapolate:} Stay within range of data

\item \textbf{Understand uncertainty:} Report confidence/prediction intervals
\end{enumerate}

\vspace{0.5cm}
\textbf{In R:}
\begin{itemize}
\item \texttt{cor.test(x, y)} for correlation
\item \texttt{lm(y \textasciitilde\ x)} for regression
\item \texttt{plot(model)} for diagnostics
\item \texttt{predict(model, newdata)} for predictions
\end{itemize}
\end{frame}

\begin{frame}
\frametitle{Questions?}
\begin{center}
\Large{Thank you!}

\vspace{1cm}

\textbf{Email:} john.molitor@oregonstate.edu\\
\textbf{Office:} 153 Milam
\end{center}
\end{frame}

\end{document}
